\chapter{The Empirics of Price Impact}\label{ch:mx:the-empirics-of-price-impact}

Buy one per cent of a day's volume and the price moves, on average, by some amount; buy four per cent and it moves about twice as much, not four times.
That the impact of an order grows like the square root of its size is one of the most robust regularities in finance, and nobody designed it. It is also
hard to measure: impact is a small average under a large noise, the orders that are measured were sent for reasons that also move prices, and a simulated
market can show impact only if it can respond. This chapter measures impact at three scales, a single trade, the flow of an interval and a metaorder, in
a reactive simulated market, then fits the square-root law on the firm's metaorder records and shows how a forecast in their timing biases it.

\section{The impact of a single trade}

Impact is the price move that follows a trade in its direction (One Quant Book 7, chapter~18); chapter~5 measured it per trade, against the mid, at a
horizon. The chapter's market is \texttt{firm.agentmkt}'s population on \texttt{firm.exchsim}: liquidity providers who place one-lot orders relative to
the book as it is, a share one tick from the opposite quote and the rest up to ten ticks away with more weight further out (a book whose depth grows away
from the price), noise takers whose orders arrive at pre-drawn times with pre-drawn signs and sizes, and fundamentalists who trade toward a hidden value
that moves by random jumps. Everything but the noise orders and the value's jumps reacts to the book.

\begin{definition}[Impact curve]\label{def:mx:the-empirics-of-price-impact:curve}
An \emph{impact curve}\index{impact curve} is the mean price move in the direction of a trade, or of a metaorder, as a function of the time since it
started (the response function of a single trade, in trades or seconds) or of its size.
\end{definition}

The response to a single \omterm{def:mx:the-limit-order-book:orders}{market order}, merged from its executions, grows with its size in the simulated hour (763\,900 shares in 7\,639 trades): after one
more order it is 0.27 ticks for one lot, 0.58 for two or three, 1.12 for four to seven and 2.34 for eight or more (\cref{fig:mx:the-empirics-of-price-impact:single},
left). It then decays for large orders, to 0.46 ticks fifty orders later, while it grows for one-lot orders, to 0.81: the fundamentalists who know the
value send one-lot orders, and their trades predict the price. In this thin book the response is close to proportional to size; in real markets it is
strongly concave, partly because large \omterm{def:mx:the-limit-order-book:orders}{market orders} are sent when the book can take them.

\begin{definition}[Aggregate impact]\label{def:mx:the-empirics-of-price-impact:aggregate}
\emph{Aggregate impact}\index{aggregate impact} is the mean price change over an interval as a function of the interval's signed order flow (buys less
sells), whoever traded.
\end{definition}

Over ten-second intervals the mid moves 2.1 ticks per thousand shares of net buying; the binned means run from $-3.1$ ticks for about $-1\,400$ shares to
$+3.1$ for $+1\,300$ (\cref{fig:mx:the-empirics-of-price-impact:single}, right). \omterm{def:mx:the-empirics-of-price-impact:aggregate}{Aggregate impact} is linear for small imbalances and bends at large ones,
where the flow meets the book's deeper levels and the fundamentalists.

\begin{omfigure}
\begin{tikzpicture}
\begin{groupplot}[group style={group size=2 by 1,horizontal sep=1.6cm},width=6.6cm,height=5cm,
  tick label style={font=\footnotesize},label style={font=\small},grid=major,grid style={gray!25},table/col sep=comma]
\nextgroupplot[title={\small response to one order (ticks)},xmode=log,xlabel={orders later},ymin=0,ymax=2.6,
  legend columns=2,legend style={font=\footnotesize,at={(0.5,-0.3)},anchor=north,draw=none,fill=none}]
\addplot[omDef,very thick,mark=*] table[x=lag,y=s0]{figdata/microstructure/11-the-empirics-of-price-impact/response.csv};
\addplot[omThm,very thick,mark=square*] table[x=lag,y=s1]{figdata/microstructure/11-the-empirics-of-price-impact/response.csv};
\addplot[omProp,very thick,mark=triangle*] table[x=lag,y=s2]{figdata/microstructure/11-the-empirics-of-price-impact/response.csv};
\addplot[black,very thick,mark=diamond*] table[x=lag,y=s3]{figdata/microstructure/11-the-empirics-of-price-impact/response.csv};
\legend{1 lot,{2--3 lots},{4--7 lots},{8+ lots}}
\nextgroupplot[title={\small 10-second mid change (ticks)},xlabel={signed volume (shares)},xtick={-2000,0,2000},
  scaled x ticks=false,xticklabel style={/pgf/number format/.cd,fixed,1000 sep={\,}}]
\addplot[omDef,very thick,mark=*] table[x=flow,y=dm]{figdata/microstructure/11-the-empirics-of-price-impact/aggregate.csv};
\end{groupplot}
\end{tikzpicture}

\omcaption{Impact at two scales in the simulated market. Left: the mean mid move in the direction of a \omterm{def:mx:the-limit-order-book:orders}{market order}, by its size, against the number of
orders since. Right: the mid change over ten-second intervals against their net signed volume, binned. Data: \texttt{mx\_impact.single\_and\_aggregate}.}\label{fig:mx:the-empirics-of-price-impact:single}
\end{omfigure}

\section{The impact of a metaorder: during and after}

\begin{definition}[Peak impact and impact reversion]\label{def:mx:the-empirics-of-price-impact:peak}
The \emph{peak impact}\index{peak impact} of a metaorder is the price move in its direction from its start to its last fill. \emph{Impact
reversion}\index{impact reversion} is the part of that move that disappears after the metaorder ends; what remains after a long time is its permanent
impact.
\end{definition}

A metaorder's impact is a few ticks under a price that wanders by several ticks a minute, so a single execution says little. A simulator can do better
than the data: run the session twice, with and without the metaorder, with the same pre-drawn noise orders and value jumps, and take the difference of
the two price paths. The noise that does not depend on the metaorder cancels; what the providers and fundamentalists do differently does not, because a
small change in the book changes their later choices.

Two results follow. First, a replayed background cannot show impact at all. Sent into \texttt{firm.tape}'s replayed flow, a 24\,000-share metaorder
leaves the mid path exactly where it was without it, at every horizon from 5 to 20 minutes: the replayed liquidity providers keep posting at the prices
they posted at in the recording and never learn of the order. Second, in the reactive market (\cref{fig:mx:the-empirics-of-price-impact:meta}), 30 children
over ten minutes of a 24\,000-share buy, 3.2\% of the hour's 749\,200 shares, move the price 6.3 ticks halfway through (standard error 1.9 over six sessions)
and 7.8 (2.2) at the end; five minutes after the end it is $-1.5$ (1.5): the fundamentalists have pulled it back. The 6\,000- and 1\,500-share orders end at
2.6 (1.9) and 2.4 (0.9) ticks, within their errors of each other. Six sessions resolve the largest order and not the curve; the square-root law was
found in hundreds of thousands of real metaorders, as in Bacry, Iuga, Lasnier and Lehalle (2015), and in the trading profiles of Moro and co-authors (2009).

\begin{omfigure}
\begin{tikzpicture}
\begin{axis}[width=9.5cm,height=5.2cm,xlabel={\small seconds from the start of the metaorder},ylabel={\small impact (ticks)},
  tick label style={font=\footnotesize},xmin=-50,xmax=1250,ymin=-6,ymax=12,grid=major,grid style={gray!25},
  xticklabel style={/pgf/number format/.cd,fixed,1000 sep={\,}},
  legend columns=3,legend style={font=\footnotesize,at={(0.5,-0.25)},anchor=north,draw=none,fill=none,
  /tikz/every even column/.append style={column sep=8pt}},table/col sep=comma]
\fill[gray!12] (axis cs:0,-6) rectangle (axis cs:600,12);
\addplot[omDef,thick,mark=*,error bars/.cd,y dir=both,y explicit] table[x=h,y=m1500,y error=e1500]{figdata/microstructure/11-the-empirics-of-price-impact/meta.csv};
\addplot[omThm,thick,mark=square*,error bars/.cd,y dir=both,y explicit] table[x=h,y=m6000,y error=e6000]{figdata/microstructure/11-the-empirics-of-price-impact/meta.csv};
\addplot[black,very thick,mark=diamond*,error bars/.cd,y dir=both,y explicit] table[x=h,y=m24000,y error=e24000]{figdata/microstructure/11-the-empirics-of-price-impact/meta.csv};
\legend{1\,500 shares,6\,000 shares,24\,000 shares}
\end{axis}
\end{tikzpicture}

\omcaption{Counterfactual impact of metaorders in the simulated market: each session with the metaorder less the same session without it, mean and
standard error over six sessions; the shaded band is the execution. Data: \texttt{mx\_impact.counterfactual}.}\label{fig:mx:the-empirics-of-price-impact:meta}
\end{omfigure}

\section{The square-root law}

\begin{definition}[Impact prefactor]\label{def:mx:the-empirics-of-price-impact:prefactor}
In the square-root law $I=Y\sigma\,(Q/V)^{\delta}$ with $\delta\approx\tfrac12$, where $I$ is the \omterm{def:mx:the-empirics-of-price-impact:peak}{peak impact} as a fraction of the price, $\sigma$ the daily
volatility, $Q$ the metaorder's size and $V$ the daily volume, the \emph{impact prefactor}\index{impact prefactor} $Y$ is the dimensionless constant, of
order one, that sets the level.
\end{definition}

With $Y=0.8$ and $\sigma=2\%$, buying 1\% of the day's volume costs about $0.8\times0.02\times\sqrt{0.01}=16$ basis points of impact and 4\% about 32: four
times the size, twice the impact. Toth and co-authors (2011) explained the law by a book whose average supply and demand are V-shaped and vanish at the
price, so that large metaorders must be split and the local liquidity is anomalously small; Zarinelli, Treccani, Farmer and Lillo (2015) found that in US
data the square root fits over about two orders of magnitude of size and a logarithm over almost five. Loeb (1983) had tabulated trading costs rising
with order size long before.

The firm's records are the right place to fit it: many metaorders across stocks, with their sizes, durations, volatilities and fills. The chapter
simulates them: 2\,000 metaorders on 50 stocks, sizes from 0.01\% to 5\% of the daily volume, durations at a tenth of the volume, a true law $Y=0.8$,
$\delta=\tfrac12$, and a price noise of $\sigma\sqrt T$ on each. \texttt{firm.impactfit.fit\_power} averages impact over sizes within bins (a single
impact can be negative, so logs of individual impacts are useless), fits the log of the bin means on the log of the size (Book~7's
\texttt{firm.tcost.fit\_impact}), and bootstraps whole stocks for the intervals.

\omcode{code/firm/impactfit/firm_impactfit.py}{83}{103}{The exponent and prefactor of the square-root law, with 95\% intervals from resampling whole stocks.}\label{lst:mx:the-empirics-of-price-impact:fit}

The fit returns $\delta=0.48$ (95\% interval 0.39 to 0.56) and $Y=0.85$ (0.55 to 1.66); with the exponent fixed at one half, $Y=0.81$. The data identify the
exponent to about $\pm0.09$ and the prefactor to a factor of two, from 2\,000 orders with the true law inside.

\section{Measurement pitfalls}

\begin{definition}[Alpha contamination]\label{def:mx:the-empirics-of-price-impact:alpha}
\emph{Alpha contamination}\index{alpha contamination} is the part of a measured impact that is the price move the order's owner expected anyway: orders
sent because of a forecast move with the forecast, and their impact looks larger than it is.
\end{definition}

Suppose 30\% of the firm's metaorders were started because a forecast expected a drift of $0.5\sigma\sqrt T$ in their direction over their duration. The
same fit returns $\delta=0.50$ (0.44 to 0.55) and $Y=1.42$ (1.01 to 2.04); with the exponent fixed, $Y=1.31$, 62\% too high
(\cref{fig:mx:the-empirics-of-price-impact:sqrt}). The exponent is untouched because the forecast's horizon is the order's duration, which grows with its
size, so the drift scales like the impact; only the level is wrong. Removing the forecast from each timed order's measured impact restores the clean fit
exactly; a regression of impact on the forecast would have found a coefficient of 1.54, not one, because larger orders carry both more impact and more
forecast.

\omcode{code/microstructure/11-the-empirics-of-price-impact/python/mx_impact.py}{94}{112}{The firm's metaorders, simulated: a square-root law, price noise, and a share of orders timed by a forecast.}\label{lst:mx:the-empirics-of-price-impact:panel}

\begin{omfigure}
\begin{tikzpicture}
\begin{axis}[width=9cm,height=5.4cm,xmode=log,ymode=log,xlabel={\small $Q/V$},ylabel={\small impact $/\,\sigma$},
  tick label style={font=\footnotesize},xmin=8e-5,xmax=0.06,ymin=0.0008,ymax=0.5,grid=major,grid style={gray!25},
  legend columns=2,legend style={font=\footnotesize,at={(0.5,-0.3)},anchor=north,draw=none,fill=none,
  /tikz/every even column/.append style={column sep=8pt}},table/col sep=comma]
\addplot[omDef,only marks,mark=*] table[x=x,y=clean]{figdata/microstructure/11-the-empirics-of-price-impact/sqrt.csv};
\addplot[omThm,only marks,mark=square*] table[x=x,y=timed]{figdata/microstructure/11-the-empirics-of-price-impact/sqrt.csv};
\addplot[omDef,domain=1e-4:0.05,samples=20] {0.81*sqrt(x)};
\addplot[omThm,dashed,domain=1e-4:0.05,samples=20] {1.31*sqrt(x)};
\legend{forecast removed,with timed orders,$0.81\sqrt{Q/V}$,$1.31\sqrt{Q/V}$}
\end{axis}
\end{tikzpicture}

\omcaption{The square-root law on 2,000 simulated metaorders: binned mean impact over daily volatility against size over daily volume, with the
forecast removed and with 30\% of the orders timed by a forecast, and the square-root fits. The slope is the same; the level is not. Data:
\texttt{mx\_impact.panel\_fits}.}\label{fig:mx:the-empirics-of-price-impact:sqrt}
\end{omfigure}

The other pitfalls are as common. Noise: the price moves $\sigma\sqrt T$ over an order's life whatever the order does, so small orders' impact is
invisible without many of them, and measuring it over a long fixed horizon biases the exponent down. Selection: only executed orders are recorded, and
orders that were cancelled because the price ran away are missing. Overlap: concurrent metaorders of the same firm, or of the market, add their impact to
each other's. Replay: a background that cannot respond shows no impact at all. Every published estimate is a choice about these, which is why
estimates of $Y$ differ more than estimates of $\delta$.

\section{Tutorial: the square root in our own data}\label{tut:mx:the-empirics-of-price-impact:tut}

\begin{tutorial}
\textbf{Goal.} Measure impact at three scales in a reactive market, then fit the square-root law on the firm's metaorders and correct it for their timing.
\textbf{End state:} \cref{fig:mx:the-empirics-of-price-impact:single,fig:mx:the-empirics-of-price-impact:meta,fig:mx:the-empirics-of-price-impact:sqrt}.
\begin{tutsteps}
\item \textbf{Market.} \texttt{firm\_agentmkt.session(mx\_impact.MARKET, seconds, seed, agents)}; \texttt{res.tape()} for trades and quotes.
\item \textbf{Single and aggregate.} \texttt{firm\_impactfit.response(trades, times, mids)} and \texttt{aggregate(\dots, 10.0, edges)}.
\item \textbf{Metaorders.} \texttt{mx\_impact.counterfactual()}: each session with and without a \texttt{Metaorder}; \texttt{replay\_pitfall()} for the
replayed flow.
\item \textbf{Law.} \texttt{panel(seed, alpha\_share)} and \texttt{fit\_power(impact, sigma, participation, clusters=stock)};
\texttt{decontaminate(impact, forecast)}; draw with \texttt{fig\_impact.py}.
\end{tutsteps}
\textbf{What to change next.} Give the providers a flatter placement law (shape 0) and see whether the single-order response becomes concave; fit the
panel with a logarithm instead of a power and compare the residuals.
\end{tutorial}

\section{Build: impact estimation}\label{bld:mx:the-empirics-of-price-impact:impactfit}

\begin{build}
\bfield{Purpose} \omterm{def:mx:the-empirics-of-price-impact:curve}{Impact curves} and laws from the firm's fills and market data, for the execution models of chapters 14 and 15, the pre-trade estimates of
chapter~19 and Book~7's cost model.
\bfield{Interface} \texttt{orders(trades)}, \texttt{response(trades, times, mids, lags, size\_edges)}, \texttt{aggregate(trades, times, mids, interval,
edges)}, \texttt{metaorder\_path(times, mids, start, end, sign, grid)}, \texttt{fit\_power(impact, sigma, participation, bins, boot, clusters)},
\texttt{decontaminate(impact, forecast)}. With it, \texttt{firm.agentmkt} (\texttt{PopulationConfig}, \texttt{Population}, \texttt{Metaorder},
\texttt{session}), whose full description is chapter~27's.
\bfield{Rules} Executions at one time and sign are one order; responses in orders; binned fits (Book~7's \texttt{fit\_impact}); cluster bootstrap by
stock; forecasts subtracted order by order.
\bfield{Acceptance tests} \texttt{code/firm/impactfit/tests/}: merged orders and a response by hand; the slope of a linear \omterm{def:mx:the-empirics-of-price-impact:aggregate}{aggregate impact}; the exponent
and prefactor of a planted law inside their intervals. \texttt{code/firm/agentmkt/tests/}: a deterministic population, shared exogenous flow, persistent
signs, a metaorder that trades its plan.
\bfield{Stretch} Impact surfaces in size and duration; decay kernels after completion (chapter~12); the logarithmic law.
\end{build}

\begin{omsources}
\item T.~F.~Loeb, ``Trading cost: the critical link between investment information and results'', \emph{Financial Analysts Journal} 39(3), 1983.
\item E.~Moro, J.~Vicente, L.~G.~Moyano, A.~Gerig, J.~D.~Farmer, G.~Vaglica, F.~Lillo and R.~N.~Mantegna, ``Market impact and trading profile of hidden
orders in stock markets'', \emph{Physical Review E} 80, 2009.
\item B.~T\'oth, Y.~Lemperi\`ere, C.~Deremble, J.~de~Lataillade, J.~Kockelkoren and J.-P.~Bouchaud, ``Anomalous price impact and the critical nature of
liquidity in financial markets'', \emph{Physical Review X} 1, 2011.
\item E.~Bacry, A.~Iuga, M.~Lasnier and C.-A.~Lehalle, ``Market impacts and the life cycle of investors orders'', \emph{Market Microstructure and
Liquidity} 1(2), 2015.
\item E.~Zarinelli, M.~Treccani, J.~D.~Farmer and F.~Lillo, ``Beyond the square root: evidence for logarithmic dependence of market impact on size and
participation rate'', \emph{Market Microstructure and Liquidity} 1(2), 2015.
\end{omsources}

\section{Exercises}

\begin{exercise}[$\star$]\label{exo:mx:the-empirics-of-price-impact:1}
With $Y=0.8$, $\delta=\tfrac12$ and a daily volatility of 2\%, what impact does the square-root law give for buying 1\% and 4\% of the day's volume?
\end{exercise}

\begin{exercise}[$\star$]\label{exo:mx:the-empirics-of-price-impact:2}
In the simulated market, the mid moves 2.1 ticks per thousand shares of net buying over ten seconds. How far should 600 shares of net buying move it, and
why is the relation only approximately linear?
\end{exercise}

\begin{exercise}[$\star$]\label{exo:mx:the-empirics-of-price-impact:3}
The 24\,000-share metaorder moved the price 7.8 ticks by its end and $-1.5$ five minutes later. How much of its impact reverted, and what in the market
caused it?
\end{exercise}

\begin{exercise}[$\star\star$]\label{exo:mx:the-empirics-of-price-impact:4}
Why does a forecast whose horizon is the order's duration bias the prefactor but not the exponent?
\end{exercise}

\begin{exercise}[$\star\star$]\label{exo:mx:the-empirics-of-price-impact:5}
Why are the intervals computed by resampling whole stocks rather than single orders?
\end{exercise}

\begin{exercise}[$\star\star$]\label{exo:mx:the-empirics-of-price-impact:6}
Why does a replayed background show no impact, and what does that imply for backtests on replayed data?
\end{exercise}

\begin{exercise}[$\star\star\star$]\label{exo:mx:the-empirics-of-price-impact:7}
\emph{Coding.} Rerun the panel with every order's duration fixed at a tenth of a day, with and without the timed orders. Fit the law and explain what you
find.
\end{exercise}

\begin{exercise}[$\star\star\star$]\label{exo:mx:the-empirics-of-price-impact:8}
\emph{Find the flaw.} ``Our fills give an exponent of 0.5, exactly the published value, so our impact model is right and our alpha is not in it.''
\end{exercise}

\section{Problem: The Square Root in Our Own Data}

\begin{problem}[{Weekend problem --- the square root in our own data}]\label{pb:mx:the-empirics-of-price-impact:1}
The execution desk wants a pre-trade impact model and asks research to fit it on the firm's own metaorders. Measure, fit and check it.

\medskip\noindent\textbf{Part I --- Small scales.}
\begin{enumerate}
\item Describe the simulated market and what in it reacts.
\item Give the response to a \omterm{def:mx:the-limit-order-book:orders}{market order} by size after one order, and after fifty.
\item Why does the response of one-lot orders grow with the lag?
\item Give the \omterm{def:mx:the-empirics-of-price-impact:aggregate}{aggregate impact} per thousand shares over ten seconds.
\end{enumerate}

\medskip\noindent\textbf{Part II --- Metaorders.}
\begin{enumerate}[resume]
\item How is a counterfactual impact measured in a simulator, and what does it remove?
\item What does the replayed background show, and why?
\item Give the 24,000-share order's impact halfway, at the end and five minutes after.
\item Why do six sessions not resolve the \omterm{def:mx:the-empirics-of-price-impact:curve}{impact curve}?
\end{enumerate}

\medskip\noindent\textbf{Part III --- The law.}
\begin{enumerate}[resume]
\item State the square-root law and compute exercise~1.
\item How is the law fitted on the panel, and why in bins?
\item Give the exponent and prefactor with their intervals.
\item What did Toth and co-authors, and Zarinelli and co-authors, find?
\end{enumerate}

\medskip\noindent\textbf{Part IV --- Pitfalls and the verdict.}
\begin{enumerate}[resume]
\item Define \omterm{def:mx:the-empirics-of-price-impact:alpha}{alpha contamination}.
\item Give the fit with 30\% of orders timed by a forecast.
\item Why is the exponent unbiased here?
\item How is the forecast removed, and what does a regression on it return?
\item What happens with a fixed long horizon (exercise~7)?
\item List three other measurement pitfalls.
\item State the \emph{named result}: the exponent and prefactor of the square-root law fitted on the firm's simulated metaorders with confidence intervals,
and their bias when order timing carries a forecast.
\item In one sentence: what should the desk's pre-trade model use?
\end{enumerate}
\end{problem}

\section{Interview questions}

\begin{interviewq}[$\star$ \iqroles{trader, researcher}]\label{iq:mx:the-empirics-of-price-impact:1}
State the square-root impact law and what each term means.
\end{interviewq}

\begin{interviewq}[$\star\star$ \iqroles{researcher}]\label{iq:mx:the-empirics-of-price-impact:2}
How would you estimate the impact of your firm's metaorders, and what would bias the estimate?
\end{interviewq}

\begin{interviewq}[$\star\star$ \iqroles{researcher}]\label{iq:mx:the-empirics-of-price-impact:3}
Why is impact concave in size? Give one explanation.
\end{interviewq}

\begin{interviewq}[$\star\star$ \iqroles{trader}]\label{iq:mx:the-empirics-of-price-impact:4}
Your backtest on replayed market data shows no cost for a large order beyond the spread. What is wrong?
\end{interviewq}

\begin{interviewq}[$\star\star$ \iqroles{researcher}]\label{iq:mx:the-empirics-of-price-impact:5}
How much of a metaorder's impact is permanent, and what does the answer depend on?
\end{interviewq}

\begin{interviewq}[$\star\star\star$ \iqroles{researcher}]\label{iq:mx:the-empirics-of-price-impact:6}
Your alpha team's orders show twice the impact of the index team's orders of the same size. Is their execution worse?
\end{interviewq}
